%   Filename    : abstract.tex

\begin{center}
\textbf{Abstract}
\end{center}

\setlength{\parindent}{0pt}

Next-generation sequencing (NGS) platforms have advanced genomic research but remain susceptible to PCR-induced chimeras that can compromise mitochondrial genome assembly. These artificial hybrid reads are especially problematic for small, circular mitochondrial genomes, where they may produce fragmented contigs and false junctions. Existing chimera detection tools are mainly designed for amplicon-based microbial analysis and often rely on reference databases or abundance assumptions that are unsuitable for organellar assembly. To address this gap, this study presents MitoChime, a machine learning pipeline for detecting PCR-induced chimeric reads in \textit{Sardinella lemuru} Illumina paired-end data.

Simulated clean and chimeric reads were generated from a mitochondrial reference genome, aligned, and used to construct labelled datasets for feature-based and sequence-based modelling. Alignment-derived and sequence-derived features were extracted to benchmark classical machine learning classifiers, while a one-dimensional convolutional neural network (CNN) and a Bidirectional Gated Recurrent Unit (BiGRU) were developed for direct sequence-based classification. Among the classical models, tuned Gradient Boosting achieved the strongest held-out performance with an F1-score of 0.7765 and ROC-AUC of 0.8459. Feature analysis showed that clipping-based and k-mer divergence features contributed most of the predictive signal, while microhomology features added minimal value. The deep learning models performed substantially better, with the CNN achieving an F1-score of 0.9135 and ROC-AUC of 0.9607, and the BiGRU achieving the best overall performance with an F1-score of 0.9523 and ROC-AUC of 0.9849. External assembly testing showed that filtering improved assembly structure, although the best read-level classifier, BiGRU, did not always produce a continuous assembly graph. These findings indicate that MitoChime is a useful pre-assembly step for mitochondrial Illumina workflows, and that model selection should consider both predictive performance and downstream assembly effects.

\begin{tabular}{lp{4.25in}}
\hspace{-0.5em}\textbf{Keywords:}\hspace{0.25em} & Chimera detection, mitochondrial genome, assembly, machine learning
\end{tabular}